
\documentclass{article}
\usepackage{colortbl}
\usepackage{makecell}
\usepackage{multirow}
\usepackage{supertabular}

\begin{document}

\newcounter{utterance}

\centering \large Interaction Transcript for game `ifeval', experiment `ifeval', episode 54 with Qwen3.6{-}27B.
\vspace{24pt}

{ \footnotesize  \setcounter{utterance}{1}
\setlength{\tabcolsep}{0pt}
\begin{supertabular}{c@{$\;$}|p{.15\linewidth}@{}p{.15\linewidth}p{.15\linewidth}p{.15\linewidth}p{.15\linewidth}p{.15\linewidth}}
    \# & \multicolumn{2}{c}{Player} && \multicolumn{2}{c}{Game Master} \\
    \hline

    \theutterance \stepcounter{utterance}  
    & & & \multicolumn{4}{p{0.6\linewidth}}{
        \cellcolor[rgb]{0.9,0.9,0.9}{
            \makecell[{{p{\linewidth}}}]{
                \texttt{\tiny{[P1$\langle$GM]}}
                \texttt{Take the text below as a starting point, and make it a complete article: "You may have to meet with a helper to work out a parenting plan. The first would be to estimate how many time you have everyday for parenting, and is that enough...."} \\
\texttt{Avoid using the following keywords: sleep, cook, feed} \\
\texttt{Mention the keyword 'schedule' for more than 5 times.} \\
            }
        }
    }
     \\ \\

    \theutterance \stepcounter{utterance}  
    & \multicolumn{4}{p{0.6\linewidth}}{
        \cellcolor[rgb]{0.561,0.737,0.561}{
            \makecell[{{p{\linewidth}}}]{
                \texttt{\tiny{[P1$\rangle$GM]}}
                \texttt{\# Navigating the Path to a Balanced Parenting Plan} \\
\\ 
\texttt{Creating a functional parenting plan is often one of the most challenging yet necessary steps for co{-}parents. You may have to meet with a helper, such as a mediator or family counselor, to work out a parenting plan that serves the best interests of your children. The first step in this process is to estimate how much time you have every day for parenting and determine if that amount is sufficient to meet your child’s emotional and developmental needs.} \\
\\ 
\texttt{This initial assessment requires honesty and precision. It is not merely about counting hours but understanding the quality of those hours. A well{-}structured **schedule** provides the framework within which these interactions occur, ensuring consistency and predictability for the child. Without a clear **schedule**, days can become chaotic, leading to stress for both parents and children alike.} \\
\\ 
\texttt{When you sit down with a helper, they will likely guide you through creating a detailed **schedule** that accounts for school hours, extracurricular activities, and necessary downtime. This **schedule** must be realistic, acknowledging the logistical constraints of work and transportation. It is crucial to review this **schedule** regularly, as children’s needs evolve, and what works in one season may not work in another.} \\
\\ 
\texttt{Furthermore, the **schedule** should include specific provisions for communication between parents. Clear guidelines on how and when to discuss changes to the **schedule** can prevent misunderstandings and conflict. By establishing a robust **schedule**, you create a stable environment where your child feels secure and supported.} \\
\\ 
\texttt{Ultimately, the goal is to ensure that the time spent parenting is meaningful and effective. Whether you are adjusting the **schedule** for holidays, school breaks, or daily routines, the focus must remain on the child’s well{-}being. A thoughtful, adaptable **schedule** is the cornerstone of a successful co{-}parenting relationship, allowing both parents to fulfill their roles with clarity and confidence.} \\
            }
        }
    }
    & & \\ \\

    \theutterance \stepcounter{utterance}  
    & & & \multicolumn{2}{p{0.3\linewidth}}{
        \cellcolor[rgb]{0.9,0.9,0.9}{
            \makecell[{{p{\linewidth}}}]{
                \texttt{\tiny{[GM$|$GM]}}
                {'keywords:forbidden_words': {'forbidden_words': ['sleep', 'cook', 'feed']}, 'keywords:frequency': {'relation': 'at least', 'keyword': 'schedule', 'frequency': 6}}
            }
        }
    }
    & & \\ \\

    \theutterance \stepcounter{utterance}  
    & & & \multicolumn{2}{p{0.3\linewidth}}{
        \cellcolor[rgb]{0.9,0.9,0.9}{
            \makecell[{{p{\linewidth}}}]{
                \texttt{\tiny{[GM$|$GM]}}
                \texttt{game\_result = WIN} \\
            }
        }
    }
    & & \\ \\

\end{supertabular}
}

\end{document}
